% ============================================================================
% Chapter 2 — Background and Related Work
% ============================================================================
\chapter{Background and Related Work}
\label{ch:background}

This chapter situates \versimux{} within the broader landscape of automated usability evaluation, simulated user testing, multi-agent systems, and automated program repair. Each section reviews the state of the art, identifies limitations that motivate the design of \versimux{}, and positions the contributions of this thesis relative to existing work.


% ---- 2.1 Automated Usability and Accessibility Evaluation --------------------
\section{Automated Usability and Accessibility Evaluation}
\label{sec:bg-accessibility}

Web accessibility evaluation, in its automated form, rests on a deceptively simple premise: that a meaningful subset of the \acr{WCAG} success criteria can be reduced to deterministic predicates over \acr{DOM} structure and computed styles. The \acr{WCAG} 2.1 standard~\cite{wcag21}, extended by version 2.2~\cite{wcag22}, organises its requirements around four principles---Perceivable, Operable, Understandable, and Robust---spanning 78 success criteria across three conformance levels (A, AA, AAA). Of these criteria, roughly 30--40\% admit fully automated verification: whether an image element carries alternative text, whether interactive controls are keyboard-reachable, whether colour contrast ratios exceed specified thresholds~\cite{abascal2004guidelines}. The remainder---criteria involving semantic clarity, cognitive load, error recovery, and consistent navigation---require human judgement that no static rule can replicate.

The most widely adopted automated checker is axe-core~\cite{axecore}, an open-source JavaScript engine developed by Deque Systems that implements over 100 \acr{WCAG} rules as \acr{DOM} traversal predicates. When injected into a browser context, axe-core queries the accessibility tree and computed style properties to produce a structured violation list containing the offending element, the violated rule identifier, and a remediation hint. Its architecture is deliberately stateless: each invocation evaluates the page as a frozen snapshot, without performing interactions or tracking state transitions. This makes axe-core fast and deterministic---properties that have driven its integration into continuous integration pipelines, browser extensions, and quality gates worldwide---but it also confines its coverage to the subset of \acr{WCAG} criteria that can be expressed as structural predicates.

Google Lighthouse~\cite{lighthouse} extends the automated auditing paradigm beyond accessibility to encompass performance, progressive web app compliance, and search engine optimisation. Its accessibility audit draws heavily on axe-core's rule set but supplements it with additional heuristics derived from Chrome's accessibility tree implementation. Lighthouse produces a numerical score (0--100) that aggregates individual rule pass rates, weighted by severity. While this scoring approach provides a convenient summary metric, Brajnik~\cite{brajnik2008comparative} demonstrated in a comparative study of evaluation methods that significant variation in the types and quantities of issues detected by each method suggests that no single automated approach---and by extension, no single aggregated score---provides a reliable picture of accessibility conformance.

Several additional tools occupy distinct niches within the automated evaluation landscape. WAVE~\cite{wave}, developed by the WebAIM initiative, overlays visual annotations directly onto the rendered page, marking violations with colour-coded icons that allow developers to inspect issues in their visual context. Pa11y~\cite{pa11y} provides a command-line interface and Node.js \acr{API} that make it straightforward to integrate accessibility checks into continuous deployment scripts, supporting multiple reporting formats and configurable rule sets. AccessLint~\cite{accesslint} takes a different architectural approach entirely, embedding itself into the GitHub pull request workflow and posting violation comments directly on the lines of code that introduce new accessibility failures. Each of these tools optimises for a different stage of the development lifecycle, yet all share the same fundamental constraint: they evaluate code, not user experience.

The limitations of static rule checking have been recognised since the earliest efforts to automate accessibility verification. Abascal et al.~\cite{abascal2004guidelines} showed that guideline-based checkers achieve high precision for structural rules but produce both false positives---flagging technically non-compliant patterns that are functionally accessible---and false negatives---missing violations that depend on interaction context. More recently, the WebAIM Million report~\cite{webaim_million}, which annually evaluates the home pages of the top one million websites, has consistently found that pages scoring well on automated checks still exhibit pervasive usability barriers when assessed by users with disabilities. The implication is clear: automated accessibility evaluation as currently practised addresses only the most tractable portion of the problem, leaving the interaction-dependent, context-sensitive, and user-subjective dimensions of usability and accessibility largely unexamined.

\begin{table}[htbp]
  \centering
  \caption{Comparison of existing automated accessibility evaluation tools.}
  \label{tab:tool-comparison}
  \begin{tabularx}{\textwidth}{@{}l L{2.5cm} L{2.5cm} C{1.8cm} C{1.8cm} C{1.8cm}@{}}
    \toprule
    \textbf{Tool} & \textbf{Method} & \textbf{Output Type} & \textbf{Multi-User Sim.} & \textbf{Code Patches} & \textbf{Verification} \\
    \midrule
    axe-core        & Static rule checking      & Violation list       & No  & No  & No  \\
    Lighthouse      & Static + performance      & Audit report         & No  & No  & No  \\
    WAVE            & Visual overlay            & Annotated page       & No  & No  & No  \\
    Pa11y           & Automated rule runner     & Violation list       & No  & No  & No  \\
    AccessLint      & CI/CD integration         & PR comments          & No  & No  & No  \\
    \midrule
    \versimux{}     & Simulated persona agents  & Patches + report     & Yes & Yes & Yes \\
    \bottomrule
  \end{tabularx}
\end{table}

\Cref{tab:tool-comparison} summarises the capabilities and limitations of the tools surveyed in this section. The final row positions \versimux{}, which addresses the three missing capabilities---multi-user simulation, code-level patches, and automated verification---through the agent-based architecture described in \cref{ch:overview,ch:architecture}.


% ---- 2.2 Simulated Users and Agentic Usability Testing ----------------------
\section{Simulated Users and Agentic Usability Testing}
\label{sec:bg-simulated-users}

The idea of modelling user behaviour computationally predates the web itself. The \acrfull{GOMS}{Goals, Operators, Methods, and Selection rules} family of models, systematised by John and Kieras~\cite{john1996goms}, decomposes task execution into a hierarchy of goals, cognitive and motor operators (e.g., perceiving a label, pressing a key), methods for achieving goals, and selection rules for choosing among alternative methods. The simplest member of the family, the \acrfull{KLM}{Keystroke-Level Model}, predicts task completion time by summing operator durations along a predetermined action sequence. \acr{GOMS} models are deterministic, analytically tractable, and have demonstrated strong predictive validity for routine, well-practised tasks~\cite{john1996goms}. They remain a standard reference point in \acr{HCI} curricula and usability engineering practice.

John et al.~\cite{cogtool} operationalised \acr{GOMS} prediction through CogTool, a prototyping tool that allows interface designers to define storyboard-style task sequences and automatically derive \acr{KLM}-level timing predictions from the specified interactions. CogTool made cognitive modelling accessible to practitioners who lacked formal training in cognitive psychology, demonstrating that useful performance predictions could be obtained from lightweight interface mockups rather than full implementations. However, \acr{GOMS}-based approaches share two structural limitations that constrain their applicability to accessibility evaluation. First, they require the analyst to predefine the complete action sequence---they cannot generate novel interaction paths or respond to unexpected interface states. Second, they model a single, idealised user profile and cannot represent the diversity of abilities, strategies, and assistive technology configurations that characterise real user populations.

The emergence of \acr{LLM}-based autonomous agents has fundamentally altered the landscape of simulated user testing. Rather than following predefined action sequences, \acr{LLM} agents observe the current interface state, reason about their goals, and select actions through natural language understanding of the \acr{DOM}. This capability has motivated the development of realistic evaluation environments. WebArena~\cite{zhou2024webarena} provides a self-hosted web environment comprising four fully functional web applications---e-commerce, forums, code management, and content management---that serve as standardised benchmarks for measuring agent task completion rates. Mind2Web~\cite{deng2023mind2web} complements this with a large-scale dataset of over 2,000 open-ended web tasks across 137 real-world websites, annotating the ground-truth action sequences that human users follow. Together, these resources have established the empirical foundation for evaluating whether \acr{LLM} agents can navigate complex web interfaces with sufficient fidelity to serve as proxies for human users.

Building on these foundations, several research groups have explored \acr{LLM} agents specifically for usability and design evaluation. As surveyed in \cref{sec:research-questions}, UXAgent~\cite{uxagent} demonstrated that persona-grounded agents can generate qualitative feedback that approximates the outputs of human testers, UXCascade~\cite{uxcascade} introduced scalable parallel evaluation with an interactive fix-and-retest workflow, and AgentA/B~\cite{agentab} validated the principle that agent swarms can approximate human behavioural distributions in A/B testing scenarios. These systems collectively establish that \acr{LLM}-based agents represent a viable paradigm for simulated user testing at scale.

Despite these advances, a fundamental tension pervades the paradigm. The same flexibility that enables \acr{LLM} agents to generate diverse, creative interaction sequences also allows them to produce traces that are plausible but not grounded in the actual interface state. An agent may report clicking a button that does not exist in the rendered page, describe a visual layout that differs from the computed styles, or hallucinate accessibility violations based on patterns memorised during pre-training rather than on evidence observed in the current session. Without explicit mechanisms to anchor agent perception and action in the live \acr{DOM}, the diagnostic outputs of \acr{LLM}-based evaluators remain fundamentally unverifiable---a challenge that motivates the grounding architecture discussed in \cref{sec:bg-grounding}.


% ---- 2.3 Multi-Agent Systems for Software Engineering -----------------------
\section{Multi-Agent Systems for Software Engineering}
\label{sec:bg-mas}

A multi-agent system (\acr{MAS}) consists of multiple autonomous entities---\textit{agents}---that perceive their environment, maintain internal state, act upon the environment, and coordinate with one another to achieve individual or collective objectives~\cite{wooldridge2009mas}. Wooldridge identifies four properties that distinguish agents from ordinary software components: \textit{autonomy} (operating without direct external control), \textit{reactivity} (responding to environmental changes), \textit{proactivity} (taking initiative to satisfy goals), and \textit{social ability} (interacting with other agents through defined protocols). These properties make multi-agent architectures particularly suited to problems that require diverse expertise, parallel execution, and decentralised decision-making---all characteristics of comprehensive usability evaluation.

The dominant theoretical framework for deliberative agent design is the \acrfull{BDI}{Belief-Desire-Intention} architecture, introduced by Rao and Georgeff~\cite{rao1995bdi}. In a \acr{BDI} agent, \textit{beliefs} represent the agent's current knowledge of its environment, \textit{desires} encode the goals it seeks to achieve, and \textit{intentions} constitute the committed plans it is currently executing. The deliberation cycle continuously updates beliefs from environmental perception, selects desires that are consistent with current beliefs, and commits to intention plans that advance the selected desires. This framework provides a principled basis for reasoning about agent behaviour: an agent's actions can be explained and predicted through the logical relationships among its beliefs, desires, and intentions. \versimux{}'s persona agents embody a simplified variant of the \acr{BDI} cycle: beliefs are grounded in the \acr{DOM} state extracted by the Playwright engine, desires are derived from persona task goals assigned by the supervisor, and intentions are constrained by Python-enforced action guards that prevent the agent from committing to infeasible plans.

Park et al.~\cite{park2023generative} demonstrated that \acr{LLM}-powered agents can exhibit surprisingly believable human behaviour when equipped with appropriate memory and reflection mechanisms. Their \textit{Generative Agents} architecture endows each agent with a memory stream (a time-stamped record of observations and actions), a retrieval mechanism that surfaces contextually relevant memories, and a reflection process that synthesises higher-level insights from accumulated experience. In a simulated town environment, 25 such agents formed social relationships, coordinated group activities, and exhibited emergent behaviours that human evaluators judged more believable than hand-scripted character responses. This work established the empirical viability of \acr{LLM}-based agent simulation at scale and introduced the architectural pattern of separating memory, retrieval, and reflection into distinct computational modules---a pattern that \versimux{} adapts through its WorkingMemory dataclass and its decomposition of persona cognition into a six-phase pipeline of perception, planning, decision, execution, evaluation, and reflection.

The application of \acr{MAS} architectures to software engineering tasks has yielded particularly relevant precedents. MetaGPT~\cite{metagpt} introduced a multi-agent framework in which \acr{LLM} agents are assigned distinct software engineering roles---product manager, architect, engineer, tester---and collaborate through structured message passing with standardised output artefacts (requirements documents, design diagrams, code files, test reports). By encoding human workflow conventions into the agent coordination protocol, MetaGPT demonstrated that role specialisation with structured inter-agent communication substantially reduces cascading hallucination errors compared to monolithic, single-agent approaches. SWE-agent~\cite{sweagent} complemented this finding by showing that the design of the agent--computer interface itself critically affects agent performance: providing agents with purpose-built tools for navigating codebases, editing files, and running tests improved issue resolution rates on the SWE-bench benchmark far more than simply upgrading the underlying \acr{LLM}. Both findings directly inform \versimux{}'s architecture, which assigns five specialised roles (supervisor, persona, recommender, conflict resolver, and verifier) and provides each role with domain-specific tooling (browser engine, \acr{DOM} state extractor, patch applicator) rather than relying on generic \acr{LLM} prompting alone.

The orchestration of multi-agent workflows in \versimux{} relies on LangGraph~\cite{langgraph}, a framework that models agent pipelines as directed state graphs with typed state containers. Unlike simple sequential chains or reactive event loops, LangGraph's state graph model supports three capabilities that prove essential for the \versimux{} pipeline: (1)~\textit{parallel fan-out} via the \code{Send()} primitive, enabling multiple pages and personas to be processed concurrently; (2)~\textit{conditional edges} that implement branching logic such as the correction loop triggered by failed verification; and (3)~\textit{typed state annotations} using Python's \code{Annotated} type hints, which ensure that concurrent node writes to shared state (e.g., appending results from parallel branches to a shared list) are handled safely without explicit locking. This combination of parallelism, conditionality, and type safety distinguishes LangGraph from earlier \acr{MAS} coordination mechanisms that relied on message queues, blackboard architectures, or ad hoc scripting.


% ---- 2.4 LLM Grounding, Hallucination, and Trust ---------------------------
\section{LLM Grounding, Hallucination, and Trust}
\label{sec:bg-grounding}

\acr{LLM} hallucination---the generation of content that is fluent and confident but factually incorrect or inconsistent with the provided input---represents the central reliability challenge for agentic systems operating in interactive environments. Ji et al.~\cite{ji2023hallucination} provide a comprehensive taxonomy distinguishing \textit{intrinsic} hallucinations (outputs that contradict the source material) from \textit{extrinsic} hallucinations (outputs that introduce claims neither supported nor contradicted by the source). In the context of web agent systems, hallucination manifests in particularly damaging forms: agents may reference \acr{DOM} elements that do not exist, describe visual properties that differ from the computed styles, report interaction outcomes that never occurred, or fabricate accessibility violations based on patterns memorised during pre-training rather than on evidence observed in the current page. When the purpose of the agent is to produce a diagnostic evaluation, hallucinated findings are indistinguishable from genuine ones without an independent verification mechanism.

The ReAct framework~\cite{yao2023react} addressed a closely related challenge---the disconnect between reasoning and action in language models---by proposing an interleaved thought--action--observation loop. In ReAct, the model alternates between generating free-form reasoning traces (``I need to find the search bar''), executing concrete actions (``click element \#search-input''), and incorporating observations from the environment (``the page now shows a text input with placeholder `Search\ldots'\,''). This interleaving forces the model to ground each reasoning step in an observed environmental state, reducing the accumulation of unsupported inferences that occurs in pure chain-of-thought reasoning. \versimux{}'s persona agent adopts a similar interleaved structure---the six-phase cognitive pipeline alternates between perception (\acr{DOM} state extraction), reasoning (plan and decision phases), action (Playwright execution), and evaluation (outcome assessment)---but extends it with deterministic, Python-enforced constraints that ReAct's prompt-only approach cannot guarantee.

A complementary grounding strategy, explored by Schick et al.~\cite{schick2023toolformer} in Toolformer, teaches language models to recognise when their own parametric knowledge is insufficient and to delegate specific subtasks to external tools---calculators, search engines, translation services---rather than attempting to generate answers from memory alone. The central insight is that tool use provides a \textit{verifiable} grounding mechanism: the output of a calculator can be checked, the result of a search query can be traced to a source document, and the return value of an \acr{API} call can be validated against a schema. In web agent systems, this principle translates to the use of browser automation \acr{API}s as grounding instruments: rather than asking the \acr{LLM} to describe what it ``sees'' on a page, the system extracts the \acr{DOM} state programmatically and presents it to the model as a structured observation that reflects the actual page content.

The benchmarks established by WebArena~\cite{zhou2024webarena} and Mind2Web~\cite{deng2023mind2web} have quantified the severity of the grounding problem in practice. On WebArena's task suite, even the strongest \acr{LLM} agents achieve task completion rates well below 20\%, with failure analysis revealing that a substantial fraction of errors stem from hallucinated selectors, incorrect assumptions about page state, and actions that target elements not interactable in the current \acr{DOM} configuration. These findings underscore that raw \acr{LLM} capability---even with frontier models---is insufficient for reliable web interaction without additional mechanisms to verify the feasibility and correctness of proposed actions before they are executed.

An instructive parallel exists in robotics. Ahn et al.~\cite{saycan} introduced SayCan, a system that grounds language model plans in robotic affordances: before executing a natural-language instruction (``bring me a sponge from the kitchen counter''), the system queries a learned affordance function that estimates whether each candidate action is physically feasible given the robot's current state and capabilities. Actions that the language model proposes but that the robot cannot physically execute are filtered out before they reach the motor system. This \textit{affordance grounding} pattern---using an environmental feasibility check to gate \acr{LLM} output before execution---directly informs \versimux{}'s grounding architecture. In \versimux{}, the affordance function takes the form of a \acr{DOM} selector existence check: before executing any click, type, or hover action, the Playwright engine verifies that the target selector exists in the live page and is interactable, preventing hallucinated interactions from corrupting the simulation trace.

Beyond selector verification, \versimux{} implements a multi-layered grounding architecture that includes Python-maintained working memory (ensuring the \acr{LLM} cannot contradict or fabricate its own state), five deterministic action guards (preventing degenerate interaction patterns such as infinite scrolling, repeated clicking, and observation spirals), and a post-simulation trace verifier that discards traces with excessive invalid steps before they propagate to downstream analysis. The design rationale and implementation details of these mechanisms are presented in \cref{ch:architecture,ch:ai-design}.


% ---- 2.5 Automated Program Repair and Patch Synthesis -----------------------
\section{Automated Program Repair and Patch Synthesis}
\label{sec:bg-apr}

\acrfull{APR}{Automated Program Repair} is the field concerned with automatically generating source-code modifications---patches---that correct software defects. The dominant paradigm, known as \textit{generate-and-validate}, proceeds in two phases: a generation phase that proposes candidate patches through various strategies (mutation, search, semantic analysis, or machine learning), and a validation phase that evaluates each candidate against a test suite or specification to determine whether it resolves the defect without introducing regressions~\cite{legoues2012genprog}.

Le Goues et al.~\cite{legoues2012genprog} established the field with GenProg, a search-based repair technique that applies genetic programming to the abstract syntax tree of buggy programs. GenProg defines mutation operators---insert, delete, and replace---that draw repair material from elsewhere in the same program, operating under the assumption that correct code for fixing a defect often already exists in the codebase, merely in the wrong location. Candidate patches are evaluated against a test suite, and a fitness function guides the evolutionary search towards patches that pass failing tests without breaking passing ones. While GenProg demonstrated that non-trivial bugs in large C programs could be repaired automatically, subsequent analysis revealed that a significant proportion of its patches were \textit{overfitting}: they passed the test suite by deleting functionality or suppressing error paths rather than by genuinely correcting the underlying defect.

Nguyen et al.~\cite{nguyen2013semfix} addressed the overfitting problem through SemFix, a semantic repair approach that uses symbolic execution and constraint solving to synthesise patches that are correct by construction with respect to a formal specification of the expected behaviour. By reasoning about program semantics rather than manipulating syntax, SemFix produces patches that are more likely to generalise beyond the specific test inputs used for validation. However, the computational cost of symbolic execution limits the approach to relatively small programs and localised defects, making it impractical for the scale and diversity of modifications that accessibility remediation typically requires.

The advent of \acr{LLM}s has introduced a third paradigm for program repair. Xia and Zhang~\cite{xia2023apr} conducted a systematic study of \acr{LLM}-based \acr{APR}, demonstrating that pre-trained code models can generate correct patches by treating repair as a conditional generation task: given the buggy code, the test failure message, and optionally the surrounding context, the model generates a modified code fragment. Their results showed that \acr{LLM}-based approaches substantially outperformed both search-based and semantic methods on established benchmarks, benefiting from the models' capacity to learn repair patterns from the vast corpus of code changes observed during pre-training. However, the authors noted that \acr{LLM}-generated patches are susceptible to the same hallucination risks identified in \cref{sec:bg-grounding}: models may generate syntactically plausible but semantically incorrect fixes, particularly when the repair requires domain-specific knowledge not well represented in the training data.

The application of \acr{APR} to \acr{UI} accessibility repair introduces domain-specific challenges that none of the techniques reviewed above directly address. First, \acr{UI} repairs span three distinct languages (HTML, CSS, and JavaScript), each with different semantics and failure modes---a single remediation effort that adds an \code{aria-label} attribute (HTML), adjusts a contrast ratio (CSS), and installs a keyboard event handler (JavaScript) requires coordinated modifications across all three languages for a single element. Second, \acr{UI} patches interact visually and behaviourally: two patches targeting the same element may produce conflicting styles, overlapping event handlers, or redundant \acr{ARIA} attributes. Standard \acr{APR} techniques, which validate patches against functional test suites, cannot evaluate these visual and interactive side effects. Third, the ``test suite'' for accessibility repair is not a set of unit tests but the subjective experience of users with diverse abilities---a criterion that resists straightforward programmatic specification. As discussed in \cref{sec:research-questions}, ACCESS~\cite{access} represents a notable step towards \acr{LLM}-based accessibility repair, but it operates on individual elements in isolation without conflict awareness or verification. \versimux{} addresses these challenges through a typed patch taxonomy spanning seven categories, an explicit conflict detection and resolution protocol, and a verification mechanism that re-simulates user personas on the patched interface rather than relying on traditional test suites.


% ---- 2.6 Developer-Facing Accessibility Tooling -----------------------------
\section{Developer-Facing Accessibility Tooling}
\label{sec:bg-dev-tools}

Beyond the research tools surveyed in the preceding sections, a substantial ecosystem of developer-facing accessibility tooling has emerged, designed to integrate accessibility evaluation into the daily workflows of frontend engineers. These tools vary in their integration points---browser extensions, linting rules, continuous integration hooks, code review bots---but share a common output paradigm: they identify violations and provide textual guidance on how to fix them, without generating executable code changes.

The most comprehensive commercial ecosystem is maintained by Deque Systems, the developers of axe-core~\cite{axecore}. Deque's product suite includes axe DevTools (a browser extension that provides in-context violation highlighting and guided remediation flows), axe Monitor (a continuous scanning service that tracks accessibility compliance across entire web properties over time), and axe Auditor (a workflow management platform for professional audit teams). While these tools significantly reduce the effort required to \textit{identify} accessibility issues, their remediation guidance takes the form of textual suggestions---``add an \code{alt} attribute to this image'', ``ensure this form input has a visible label''---that the developer must manually interpret and implement in code.

A parallel category of tools embeds accessibility checks directly into the software development lifecycle. AccessLint~\cite{accesslint} analyses pull requests on GitHub and posts violation comments on the specific lines of code that introduce new accessibility failures, enabling developers to address regressions before they reach production. Pa11y~\cite{pa11y} provides a command-line interface that can be integrated into continuous deployment pipelines, failing builds that exceed configurable violation thresholds. Linting plugins such as \code{eslint-plugin-jsx-a11y} check component code for common accessibility anti-patterns during development, providing immediate feedback in the editor. These integrations shift accessibility evaluation \textit{left} in the development process---closer to the point of code creation---but do not fundamentally alter the output format: the developer still receives a violation report and must independently craft the fix.

The persistent gap across the entire developer tooling landscape is the absence of \textit{actionable remediation artefacts}---concrete, executable code changes that developers can review, test, and apply rather than interpret and rewrite from scratch. This gap is not merely a convenience issue; it represents a meaningful barrier to accessibility adoption in practice. The cognitive effort required to translate a diagnostic finding (``WCAG 1.3.1: Form input missing programmatic label'') into a working fix (adding an \code{aria-label} attribute, wrapping the input in a \code{<label>} element, or restructuring the form with \code{<fieldset>} and \code{<legend>}) demands accessibility expertise that many development teams lack. By generating typed, conflict-resolved, and verified code patches, \versimux{} aims to reduce this translation burden from a specialised skill requirement to a code review task that any frontend developer can evaluate.


% ---- 2.7 Summary and Research Gap -------------------------------------------
\section{Summary and Research Gap}
\label{sec:bg-gap}

The preceding sections have surveyed six domains that collectively provide the intellectual and technical foundations for \versimux{}: automated accessibility evaluation (\cref{sec:bg-accessibility}), simulated user testing (\cref{sec:bg-simulated-users}), multi-agent system architectures (\cref{sec:bg-mas}), \acr{LLM} grounding (\cref{sec:bg-grounding}), automated program repair (\cref{sec:bg-apr}), and developer-facing tooling (\cref{sec:bg-dev-tools}). Each domain contributes essential capabilities---diagnostic rule engines, cognitive user models, agent coordination protocols, grounding mechanisms, patch generation strategies, and workflow integration patterns---but none provides a complete solution to the three gaps identified in \cref{sec:motivation}.

Static checkers and developer tools can evaluate code structure but cannot simulate user experience. Cognitive models and \acr{LLM} agents can simulate user behaviour, but without grounding they cannot guarantee that their observations correspond to reality. \acr{MAS} frameworks provide the coordination infrastructure for parallel evaluation and remediation, but they are inherently domain-agnostic and require task-specific agent designs and tools. \acr{APR} techniques can generate code patches, but they have not been adapted to the multi-language, visually interactive, conflict-prone domain of \acr{UI} accessibility. Grounding mechanisms can anchor agent actions in environmental state, but they have not been integrated into usability evaluation systems in a way that preserves both diagnostic flexibility and interaction reliability.

\begin{figure}[htbp]
  \centering
  \framebox[\textwidth]{\rule{0pt}{7.5cm}\parbox{0.9\textwidth}{\centering\large\sffamily [Figure 2.1: Related Work Domains Taxonomy vs. Research Gaps]\\ \vspace{1.2em} \normalsize\normalfont Displays the taxonomy of the six related literature domains evaluated in this chapter (Usability Testing, Accessibility Checkers, Multi-Agent Systems, Grounded AI, Program Repair, and Developer Tooling). The diagram maps each domain against the three research gaps, highlighting how \versimux{} synthesises these capabilities to close all three gaps simultaneously.}}
  \caption{Mapping of related research domains to the three identified usability and accessibility automation gaps.}
  \label{fig:bg-gaps-mapping}
\end{figure}

No existing body of work---whether a single tool, a research system, or a combination of techniques from the domains reviewed---simultaneously provides (1)~diverse, grounded persona simulation through real browser interaction, (2)~autonomous generation of typed, conflict-aware code patches across HTML, CSS, and JavaScript, and (3)~closed-loop verification through re-simulation of failing personas on the patched interface. This three-fold gap constitutes the research opportunity that \versimux{} addresses. The remainder of this thesis describes the architecture (\cref{ch:overview,ch:architecture}), \acr{LLM} integration design (\cref{ch:ai-design}), and runtime infrastructure (\cref{ch:infrastructure}) of the system, followed by an empirical evaluation (\cref{ch:evaluation}) that assesses the extent to which \versimux{} closes each gap.

